\documentclass[11pt]{article}

\usepackage[margin=25mm]{geometry}

\usepackage{amsmath,amssymb,fontspec,unicode-math}

\IfFontExistsTF{DejaVu Serif}{\setmainfont{DejaVu Serif}}{\setmainfont{Latin Modern Roman}}

\IfFontExistsTF{Latin Modern Math}{\setmathfont{Latin Modern Math}}{\setmathfont{DejaVu Math TeX Gyre}}

\usepackage{xurl,hyperref}

\setlength{\parindent}{0pt}\setlength{\parskip}{6pt}

\begin{document}

\title{Listwise Preference Diffusion Optimization for User Behavior Trajectories Prediction}

\author{ELVIA preliminary AI analysis}\date{21 October 2026}\maketitle

Preliminary assessment: Major.\par

Independent scientific review remains pending. Author adoption has not been requested.

The central ELBO bridge changes the likelihood being bounded and uses incompatible endpoint measures. Separate release defects affect NDCG direction, weighted ranking loss, sample aggregation and sampler settings; no corrected benchmark is inferred.

\section*{Review scope}

Full supplied technical report read alongside the corresponding original-paper claim passages and applicable author-release source. This pass independently checks the stated reasoning and bounds; historical executions were not rerun.

\section*{Problem 1: The NDCG helper rewards worse positions within the cutoff}

Source: Section 5.1; calculate\_hit metric helper and descending top-k call

Assessment: supported; impact: minor

\subsection*{Source claim (paraphrase)}

The release transforms zero-based relevant-item position into cutoff minus position before logarithmic discounting.

\subsection*{Analysis}

At cutoff \(K=5\), the best-ranked relevant item receives \(1/\log_2 6\), while the fifth-ranked receives \(1\). Standard single-relevant-item NDCG gives the reverse. The source requests descending sorted scores, so ascending order cannot explain the transformation. Repeating these positions across a trajectory preserves the reversal under the geometric mean.

\subsection*{Finding and scope}

The release does not compute conventional NDCG. HR is unaffected by this rank transformation; effects on published relative rankings require actual predictions and table lineage.

\subsection*{Proposed correction}

Use one-based rank \(r\) and discount \(1/\log_2(r+1)\), with zero for misses; recompute saved predictions and any affected validation selection.

The described correction remains a proposal; author adoption or a corresponding erratum is not established.

\section*{Problem 2: Multiplicative logit masking does not implement Soft-ListMLE}

Source: Equation (9); model/LPDO.py, list\_preference\_loss, lines 264–272

Assessment: supported; impact: minor

\subsection*{Source claim (paraphrase)}

The formula weights exponentiated scores of previously selected items by \(\gamma\); the active release multiplies their logits by \(\gamma\).

\subsection*{Analysis}

Use candidates \(A,B\), non-repeating target \((A,B)\), and identical score rows \((\log4,0)\). The first loss is \(\log(5/4)\) for both rules. At the second position the printed loss is \(\log(1+4\gamma)\), versus released \(\log(1+4^\gamma)\). At \(\gamma=0\) these are \(0\) and \(\log2\); at \(\gamma=1\) they agree. Score-shift invariance also holds for weighted probabilities but fails for multiplicatively masked logits.

\subsection*{Finding and scope}

The active surrogate differs even in the original tied-score, non-repeating domain. This does not establish poorer empirical predictions.

\subsection*{Proposed correction}

For non-repeating targets, add \(\log m_i\) to logits, with \(\log0=-\infty\), instead of multiplying by \(m_i\). Repeated targets require an explicit normalized sequential rule.

The described correction remains a proposal; author adoption or a corresponding erratum is not established.

\section*{Problem 3: The softened expression needs an explicit sample-space convention}

Source: Equation (9) and Section 4.3 likelihood interpretation

Assessment: not\_issue; impact: minor

\subsection*{Source claim (paraphrase)}

The softened ListMLE factors are used as a likelihood, without fully specifying how repeated trajectories receive mass.

\subsection*{Analysis}

For two candidates with scores \((\log4,0)\) and \(\gamma=0.3\), the two permutation weights are \(4/11\) and \(8/43\), totaling \(260/473<1\). However a repeat-capable sequential model with numerator weight \(\gamma\) on already selected items is normalized and retains both non-repeating path values. A permutation-mass deficit alone therefore does not refute every likelihood interpretation.

\subsection*{Finding and scope}

The repeated-outcome numerator and sample space need specification. A coherent extension exists, so this is a scoped clarification rather than an established impossibility.

\section*{Problem 4: The central ELBO derivation interchanges marginal and shared-chain likelihoods}

Source: Sections 4.3–4.4, Eqs. (12)–(15); Appendix B.2 and Appendix C

Assessment: supported; impact: major

\subsection*{Source claim (paraphrase)}

The derivation replaces a sum of item-chain joint log densities with only one chain log density, then identifies a shared-chain joint bound with a sum of marginal log-likelihoods.

\subsection*{Analysis}

For a normalized chain density \(P(x)\) and normalized emissions \(r_j\), the product rule gives \[\sum_{j=1}^k\log[P(x)r_j(i_j\mid x)]=k\log P(x)+\sum_j\log r_j(i_j\mid x).\] A shared-chain joint has coefficient one, but its marginal is \(\int P\prod_jr_j\), not the product of separate marginal integrals. With \(X\sim N(0,1)\), \(r(A\mid x)=e^{-x^2/2}\), trajectory \((A,A)\) has joint negative log-likelihood \((\log3)/2\), versus summed marginal negative log-likelihood \(\log2\). The exact joint posterior \(N(0,1/3)\) makes the one-chain variational bound equal \((\log3)/2<\log2\), so it cannot bound both objectives as asserted. Appendix B.2 also uses a point-mass embedding against a continuous kernel, yielding infinite endpoint KL, and labels conditionals over different variables a KL without a common measure.

\subsection*{Finding and scope}

The claimed tight variational derivation does not establish the stated objective. The heuristic loss can remain useful; no accuracy result is inferred false.

\subsection*{Proposed correction}

Specify a normalized sequence emission and choose either a shared-chain joint ELBO or separate marginal ELBOs with the appropriate repeated chain terms. Treat deterministic embedding endpoints with a compatible reconstruction likelihood or nondegenerate variational model.

The described correction remains a proposal; author adoption or a corresponding erratum is not established.

\section*{Problem 5: Unequal evaluation batches receive equal weight}

Source: Equation (20); main.py, evaluate aggregation

Assessment: supported; impact: minor

\subsection*{Source claim (paraphrase)}

The evaluator averages batch means rather than all user-level values.

\subsection*{Analysis}

Three complete-match indicators \((0,0,1)\), partitioned into batches of sizes \((2,1)\), produce \((0+1)/2=1/2\), although the dataset fraction is \(1/3\). The same predictions in one batch give \(1/3\). Equal-size batches are a valid special case. Sequence negative log-likelihood is also averaged this way; per-item log perplexity would additionally divide sequence NLL by fixed trajectory length.

\subsection*{Finding and scope}

The released metrics can depend on evaluation partition. Corrected tables cannot be reconstructed from aggregate means alone.

\subsection*{Proposed correction}

Accumulate sample numerators and counts, or weight each batch mean by its actual size before sequence aggregation.

The described correction remains a proposal; author adoption or a corresponding erratum is not established.

\section*{Problem 6: Released loss weights and sampling settings differ from the description}

Source: Equation (16), experimental setup; model/LPDO.py, p\_losses and sample; main.py sampler construction

Assessment: supported; impact: minor

\subsection*{Source claim (paraphrase)}

The code swaps ranking/reconstruction coefficients, scales linear noise endpoints by \(1000/T\), and iterates all configured diffusion steps.

\subsection*{Analysis}

The release computes \(\lambda L_{\rm rank}+(1-\lambda)L_{\rm MSE}+L_{\rm reg}\), while Eq. (16) assigns \(\lambda\) to reconstruction. They differ unless \(\lambda=1/2\) or component values coincide. With \(T=50\) and requested endpoints \(0.0001,0.02\), the source uses \(0.002,0.4\). Its sample loop calls the denoiser at every step, with no distinct one-step inference setting on the documented path.

\subsection*{Finding and scope}

The documented release cannot be assumed to reproduce the printed objective, numerical schedule or advertised one-step cost. The historical configuration behind each timing row is unresolved.

\subsection*{Proposed correction}

Identify the experimental source/configuration and reconcile the description. A separate consistent inference schedule is needed; setting the global training-step argument to one is not a validated one-step sampler repair.

The described correction remains a proposal; author adoption or a corresponding erratum is not established.

\section*{Limits}

Preliminary AI source/argument assessment; independent human scientific adjudication remains pending.\par

No author-provided or retained verification program was executed. Published empirical results were not independently reproduced.\par

\section*{Source identity}

NeurIPS 2025 official proceedings PDF (exact retained local copy inspected)

\url{https://proceedings.neurips.cc/paper_files/paper/2025/file/e9aa0e8679691ba5404c5b0e7caac518-Paper-Conference.pdf}

Reviewed PDF SHA-256: \nolinkurl{7966763086c08f2a6b6cb826b244e92c3cf7732209de26cecb81122cc02dfddf}

The SHA-256 identifies the exact retained PDF inspected in this rapid source/argument check. Original supplied report bytes are retained locally. The original PDF retrieval timestamp is taken from matching retained metadata.

\end{document}
